% Source: README.md: GPT-OSS-20B; Configuration; Quick Start (Python requirement only); Features (local-first model boundary)
% Converted to plain LaTeX for the AgentSouriva development reference.

\section{Models and Runtime Configuration}

\subsection{Local and External Model Boundaries}

The default model and embedding backends can operate without a separately hosted model server. For larger local models, AgentSouriva exposes an OpenAI-compatible model boundary. GPT-OSS-20B can be connected through any compatible serving setup without coupling the retrieval and grounding layers to a particular inference engine.

\subsection{GPT-OSS-20B}

AgentSouriva can connect to GPT-OSS-20B through an OpenAI-compatible model server. The configured model server must implement \texttt{POST\ /v1/chat/completions}.

%\aftabnote{RESOLVED-- in the following what is the adapter? mention what it is}

%Here, the adapter is an implementation of AgentSouriva's \texttt{ModelAdapter} Python protocol: a small translation boundary between the agent's internal question-and-evidence representation and a particular answer-generation backend. The \texttt{OpenAICompatibleModel} adapter converts a question and retrieved evidence into a chat-completions request, sends it to the configured model server, and converts the returned text and chunk identifiers into AgentSouriva's \texttt{ModelReply} structure. The local \texttt{MockModel} implements the same protocol without a server. Because the agent depends on this shared interface rather than either implementation, application logic, retrieval, and grounding do not depend on a specific serving engine.

In AgentSouriva, the adapter implements the \texttt{ModelAdapter} Python protocol, which provides a common interface between the agent's internal question-and-evidence representation and different answer-generation backends. The \texttt{OpenAICompatibleModel} adapter translates questions and retrieved evidence into OpenAI-compatible chat-completions API requests, sends them to the configured model server, and converts the responses, including generated text and referenced chunk identifiers, into AgentSouriva's \texttt{ModelReply} structure. The local \texttt{MockModel} follows the same protocol but operates without a model server. By relying on this shared interface, the agent keeps its application logic, retrieval, and grounding mechanisms independent of the underlying model-serving implementation.

Without an external model server, the default mock backend provides deterministic extractive responses for local operation and testing.

\subsection{Runtime Settings and Configuration Interface}

Runtime settings use the \texttt{SOURIVA\_} prefix. The application reads configuration during startup. The following values are shown as configuration examples in the source README, not asserted as deployment defaults.

\begingroup
\small
\setlength{\tabcolsep}{3pt}
\begin{longtable}[]{@{}p{2.25in}p{1.55in}p{2.30in}@{}}
\toprule\noalign{}
Setting & Source README example & Role \\
\midrule\noalign{}
\endhead
\bottomrule\noalign{}
\endlastfoot
\texttt{SOURIVA\_DATA\_DIR} & \texttt{./artifacts} & Persisted index and runtime artifacts \\
\texttt{SOURIVA\_CHUNK\_SIZE} & \texttt{900} & Chunk-size setting \\
\texttt{SOURIVA\_CHUNK\_OVERLAP} & \texttt{120} & Overlap between text chunks \\
\texttt{SOURIVA\_RETRIEVAL\_TOP\_K} & \texttt{4} & Number of retrieved chunks \\
\texttt{SOURIVA\_MIN\_RELEVANCE} & \texttt{0.18} & Minimum evidence-relevance threshold \\
\texttt{SOURIVA\_MODEL\_BACKEND} & \texttt{openai-compatible} & Model-backend selector \\
\texttt{SOURIVA\_MODEL\_BASE\_URL} & \texttt{http://localhost:}\newline\texttt{8001/v1} & Model server API base URL \\
\texttt{SOURIVA\_MODEL\_NAME} & \texttt{openai/gpt-oss-20b} & Model identifier \\
\texttt{SOURIVA\_MODEL\_API\_KEY} & \texttt{local} & API key supplied to the model endpoint \\
\end{longtable}
\endgroup

\subsection{Runtime Environment}

The source README specifies Python 3.11 or newer. Runtime, evaluation, tracking, MCP, and training capabilities have distinct dependency requirements, as detailed in the chronological build logs.

% Optional reviewer note: \aftabnote{Write a review note here.}
